\documentclass[a4paper,12pt]{article}

% 页面设置
\usepackage[top=2.54cm, bottom=2.54cm, left=1.91cm, right=1.91cm]{geometry}

% 中文支持
\usepackage[UTF8]{ctex}
\usepackage{fontspec}

% 数学公式支持
\usepackage{amsmath}
\usepackage{amssymb}
\usepackage{amsfonts}
\usepackage{amsthm}

\usepackage{enumitem}

% 定理环境定义
\newtheorem*{pf}{Proof}
\newtheorem*{sol}{Solution}

% 图形支持
\usepackage{graphicx}
\usepackage{tikz}

% 超链接支持
\usepackage{hyperref}
\hypersetup{
    colorlinks=true,
    linkcolor=blue,
    filecolor=magenta,
    urlcolor=cyan,
}

% 行距设置
\linespread{1.25}

% 标题信息
\title{Mathematical Analysis II\\Week 4 Homework (March 24)}
\author{袁子强\hspace{1cm} 2025533009}
% \date{}

\begin{document}

\maketitle

9. 在下列各题中，求 $\dfrac{\partial^2 z}{\partial x^2}$，$\dfrac{\partial^2 z}{\partial x\partial y}$，$\dfrac{\partial^2 z}{\partial y^2}$

(2) $z = \arctan \dfrac{x+y}{1-xy}$;
\begin{sol}
    注意到 $z = \arctan\dfrac{x+y}{1-xy} = \arctan x + \arctan y$，因此
    $$
        \frac{\partial z}{\partial x} = \frac{1}{1+x^2}, \quad \frac{\partial z}{\partial y} = \frac{1}{1+y^2}.
    $$
    从而
    \begin{itemize}
        \item $\dfrac{\partial^2 z}{\partial x^2} = \dfrac{\partial}{\partial x}\left(\dfrac{1}{1+x^2}\right) = -\dfrac{2x}{(1+x^2)^2}$
        \item $\dfrac{\partial^2 z}{\partial x\partial y} = \dfrac{\partial}{\partial x}\left(\dfrac{1}{1+y^2}\right) = 0$
        \item $\dfrac{\partial^2 z}{\partial y^2} = \dfrac{\partial}{\partial y}\left(\dfrac{1}{1+y^2}\right) = -\dfrac{2y}{(1+y^2)^2}$
    \end{itemize}
\end{sol}

(4) $z = \sin^2(ax + by)$;
\begin{sol}
    将 $z$ 化简为
    $$
        z = \frac{1 - \cos(2ax + 2by)}{2},
    $$
    因此
    $$
        \frac{\partial z}{\partial x} = a\sin(2ax + 2by), \quad \frac{\partial z}{\partial y} = b\sin(2ax + 2by).
    $$
    从而
    \begin{itemize}
        \item $\dfrac{\partial^2 z}{\partial x^2} = \dfrac{\partial}{\partial x}\left(a\sin(2ax + 2by)\right) = 2a^2\cos(2ax + 2by)$
        \item $\dfrac{\partial^2 z}{\partial x\partial y} = \dfrac{\partial}{\partial x}\left(b\sin(2ax + 2by)\right) = 2ab\cos(2ax + 2by)$
        \item $\dfrac{\partial^2 z}{\partial y^2} = \dfrac{\partial}{\partial y}\left(b\sin(2ax + 2by)\right) = 2b^2\cos(2ax + 2by)$
    \end{itemize}
\end{sol}

(6) $z = \arcsin(xy)$.
\begin{sol}
    一阶偏导数为
    $$
        \frac{\partial z}{\partial x} = \frac{y}{\sqrt{1-x^2y^2}}, \quad \frac{\partial z}{\partial y} = \frac{x}{\sqrt{1-x^2y^2}}.
    $$
    因此
    \begin{itemize}
        \item $\dfrac{\partial^2 z}{\partial x^2} = \dfrac{\partial}{\partial x}\left(\dfrac{y}{\sqrt{1-x^2y^2}}\right) = y \cdot \dfrac{xy^2}{(1-x^2y^2)^{\frac{3}{2}}} = \dfrac{xy^3}{(1-x^2y^2)^{\frac{3}{2}}}$
        \item $\dfrac{\partial^2 z}{\partial x\partial y} = \dfrac{\partial}{\partial x}\left(\dfrac{x}{\sqrt{1-x^2y^2}}\right) = \dfrac{1}{(1-x^2y^2)^{\frac{3}{2}}}$
        \item $\dfrac{\partial^2 z}{\partial y^2} = \dfrac{\partial}{\partial y}\left(\dfrac{x}{\sqrt{1-x^2y^2}}\right) = x \cdot \dfrac{x^2y}{(1-x^2y^2)^{\frac{3}{2}}} = \dfrac{x^3y}{(1-x^2y^2)^{\frac{3}{2}}}$
    \end{itemize}
\end{sol}


10. 设 $u = \mathrm{e}^{xyz}$，求 $\dfrac{\partial^3 u}{\partial x\partial y\partial z}$，$\dfrac{\partial^3 u}{\partial x\partial y^2}$.
\begin{sol}
    对于 $\dfrac{\partial^3 u}{\partial x\partial y\partial z}$，有
    \begin{align*}
        \dfrac{\partial^3 u}{\partial x\partial y\partial z} & = \dfrac{\partial^2}{\partial x\partial y}(xy\mathrm{e}^{xyz})                              \\
                                                             & = \dfrac{\partial}{\partial x}(x\mathrm{e}^{xyz} + x^2yz\mathrm{e}^{xyz})                   \\
                                                             & = \mathrm{e}^{xyz} + xyz\mathrm{e}^{xyz} + 2xyz\mathrm{e}^{xyz} + x^2y^2z^2\mathrm{e}^{xyz} \\
                                                             & = (1 + 3xyz + x^2y^2z^2)\mathrm{e}^{xyz}.
    \end{align*}

    对于 $\dfrac{\partial^3 u}{\partial x\partial y^2}$，有
    \begin{align*}
        \dfrac{\partial^3 u}{\partial x\partial y^2} & = \dfrac{\partial^2}{\partial x\partial y}(xz\mathrm{e}^{xyz}) \\
                                                     & = \dfrac{\partial}{\partial x}(x^2z^2\mathrm{e}^{xyz})         \\
                                                     & = 2xz^2\mathrm{e}^{xyz} + x^2yz^3\mathrm{e}^{xyz}.
    \end{align*}
\end{sol}


11. 设 $r = \sqrt{x^2 + y^2 + z^2}$，证明当 $r \neq 0$ 时有

(1) $\dfrac{\partial^2 r}{\partial x^2} + \dfrac{\partial^2 r}{\partial y^2} + \dfrac{\partial^2 r}{\partial z^2} = \dfrac{2}{r}$;
\begin{pf}
    一阶偏导数为
    $$
        \frac{\partial r}{\partial x} = \frac{x}{r}, \quad \frac{\partial r}{\partial y} = \frac{y}{r}, \quad \frac{\partial r}{\partial z} = \frac{z}{r}.
    $$
    因此
    $$
        \frac{\partial^2 r}{\partial x^2} = \frac{\partial}{\partial x}\left(\frac{x}{r}\right) = \frac{r - x\cdot\frac{\partial r}{\partial x}}{r^2} = \frac{r - x\cdot\frac{x}{r}}{r^2} = \frac{r^2 - x^2}{r^3}.
    $$
    同理可得
    $$
        \frac{\partial^2 r}{\partial y^2} = \frac{r^2 - y^2}{r^3}, \quad \frac{\partial^2 r}{\partial z^2} = \frac{r^2 - z^2}{r^3}.
    $$
    相加得
    $$
        \frac{\partial^2 r}{\partial x^2} + \frac{\partial^2 r}{\partial y^2} + \frac{\partial^2 r}{\partial z^2} = \frac{3r^2 - (x^2+y^2+z^2)}{r^3} = \frac{3r^2 - r^2}{r^3} = \frac{2r^2}{r^3} = \frac{2}{r}.
    $$
    Q.E.D.
\end{pf}

(2) $\dfrac{\partial^2 \ln r}{\partial x^2} + \dfrac{\partial^2 \ln r}{\partial y^2} + \dfrac{\partial^2 \ln r}{\partial z^2} = \dfrac{1}{r^2}$;
\begin{pf}
    一阶偏导数为
    $$
        \frac{\partial \ln r}{\partial x} = \frac{1}{r}\cdot\frac{\partial r}{\partial x} = \frac{x}{r^2}.
    $$
    因此二阶偏导数为
    $$
        \frac{\partial^2 \ln r}{\partial x^2} = \frac{\partial}{\partial x}\left(\frac{x}{r^2}\right) = \frac{r^2 - x\cdot 2r\cdot\frac{\partial r}{\partial x}}{r^4} = \frac{r^2 - 2x^2}{r^4}.
    $$
    同理可得
    $$
        \frac{\partial^2 \ln r}{\partial y^2} = \frac{r^2 - 2y^2}{r^4}, \quad \frac{\partial^2 \ln r}{\partial z^2} = \frac{r^2 - 2z^2}{r^4}.
    $$
    相加得
    $$
        \frac{\partial^2 \ln r}{\partial x^2} + \frac{\partial^2 \ln r}{\partial y^2} + \frac{\partial^2 \ln r}{\partial z^2} = \frac{3r^2 - 2(x^2+y^2+z^2)}{r^4} = \frac{3r^2 - 2r^2}{r^4} = \frac{1}{r^2}.
    $$
    Q.E.D.
\end{pf}

(3) $\dfrac{\partial^2}{\partial x^2}\dfrac{1}{r} + \dfrac{\partial^2}{\partial y^2}\dfrac{1}{r} + \dfrac{\partial^2}{\partial z^2}\dfrac{1}{r} = 0$.
\begin{pf}
    一阶偏导数为
    $$
        \frac{\partial}{\partial x}\left(\frac{1}{r}\right) = -\frac{1}{r^2}\cdot\frac{\partial r}{\partial x} = -\frac{x}{r^3}.
    $$
    二阶偏导数为
    $$
        \frac{\partial^2}{\partial x^2}\left(\frac{1}{r}\right) = \frac{\partial}{\partial x}\left(-\frac{x}{r^3}\right) = -\frac{r^3 - x\cdot 3r^2\cdot\frac{\partial r}{\partial x}}{r^6} = -\frac{r^3 - 3rx^2}{r^6} = \frac{3x^2 - r^2}{r^5}.
    $$
    同理可得
    $$
        \frac{\partial^2}{\partial y^2}\left(\frac{1}{r}\right) = \frac{3y^2 - r^2}{r^5}, \quad \frac{\partial^2}{\partial z^2}\left(\frac{1}{r}\right) = \frac{3z^2 - r^2}{r^5}.
    $$
    相加得
    $$
        \frac{\partial^2}{\partial x^2}\frac{1}{r} + \frac{\partial^2}{\partial y^2}\frac{1}{r} + \frac{\partial^2}{\partial z^2}\frac{1}{r} = \frac{3(x^2+y^2+z^2) - 3r^2}{r^5} = \frac{3r^2 - 3r^2}{r^5} = 0.
    $$
    Q.E.D.
\end{pf}


13. 求下列函数的微分，或在给定点的微分

(2) $z = \dfrac{xy}{x^2 + y^2}$;
\begin{sol}
    一阶偏导数为
    $$
        \frac{\partial z}{\partial x} = \frac{y(x^2+y^2) - xy\cdot 2x}{(x^2+y^2)^2} = \frac{y^3 - x^2y}{(x^2+y^2)^2} = \frac{y(y^2-x^2)}{(x^2+y^2)^2}.
    $$
    同理可得
    $$
        \frac{\partial z}{\partial y} = \frac{x(x^2-y^2)}{(x^2+y^2)^2}.
    $$
    因此
    $$
        dz = \frac{\partial z}{\partial x}dx + \frac{\partial z}{\partial y}dy = \frac{y(y^2-x^2)}{(x^2+y^2)^2}dx + \frac{x(x^2-y^2)}{(x^2+y^2)^2}dy.
    $$
\end{sol}

(4) $z = \arctan \dfrac{y}{x}$;
\begin{sol}
    一阶偏导数为
    $$
        \frac{\partial z}{\partial x} = \frac{1}{1+\left(\frac{y}{x}\right)^2}\cdot\left(-\frac{y}{x^2}\right) = -\frac{y}{x^2+y^2},
    $$
    $$
        \frac{\partial z}{\partial y} = \frac{1}{1+\left(\frac{y}{x}\right)^2}\cdot\frac{1}{x} = \frac{x}{x^2+y^2}.
    $$
    因此
    $$
        dz = -\frac{y}{x^2+y^2}dx + \frac{x}{x^2+y^2}dy.
    $$
\end{sol}

(6) $z = x^4 + y^4 - 4x^2y^2$ 在点 $(0,0),(1,1)$;
\begin{sol}
    一阶偏导数为
    $$
        \frac{\partial z}{\partial x} = 4x^3 - 8xy^2, \quad \frac{\partial z}{\partial y} = 4y^3 - 8x^2y.
    $$
    在给定点处的值为
    $$
        \begin{cases}
            \left.\dfrac{\partial z}{\partial x}\right|_{(0,0)} = 0 \\
            \left.\dfrac{\partial z}{\partial y}\right|_{(0,0)} = 0
        \end{cases}, \quad
        \begin{cases}
            \left.\dfrac{\partial z}{\partial x}\right|_{(1,1)} = 4 - 8 = -4 \\
            \left.\dfrac{\partial z}{\partial y}\right|_{(1,1)} = 4 - 8 = -4
        \end{cases}.
    $$
    因此
    $$
        \begin{cases}
            dz|_{(0,0)} = 0\cdot dx + 0\cdot dy = 0 \\
            dz|_{(1,1)} = -4dx - 4dy
        \end{cases}.
    $$
\end{sol}


15. 根据可微的定义证明，函数 $f(x,y) = \sqrt{|xy|}$ 在原点处不可微。
\begin{pf}
    即证
    $$
        \Delta z = f(\Delta x,\Delta y) - f(0,0) = f_x(0,0)\Delta x + f_y(0,0)\Delta y + o(h),
    $$
    其中 $h = \sqrt{(\Delta x)^2 + (\Delta y)^2} \to 0$。

    计算偏导数：
    $$
        f_x(0,0) = \lim_{\Delta x \to 0}\frac{f(\Delta x,0) - f(0,0)}{\Delta x} = \lim_{\Delta x \to 0}\frac{0 - 0}{\Delta x} = 0,
    $$
    $$
        f_y(0,0) = \lim_{\Delta y \to 0}\frac{f(0,\Delta y) - f(0,0)}{\Delta y} = 0.
    $$
    因此
    $$
        \lim_{h \to 0}\frac{\Delta z - \big[f_x(0,0)\Delta x + f_y(0,0)\Delta y\big]}{h} = \lim_{h \to 0}\frac{f(\Delta x,\Delta y) - f(0,0)}{h} = \lim_{\substack{\Delta x \to 0 \\ \Delta y \to 0}}\frac{\sqrt{|\Delta x\Delta y|}}{\sqrt{(\Delta x)^2 + (\Delta y)^2}}.
    $$
    取 $\Delta y = \Delta x$，则
    $$
        \lim_{\Delta x \to 0}\frac{\sqrt{|\Delta x\Delta y|}}{\sqrt{(\Delta x)^2 + (\Delta y)^2}} = \lim_{\Delta x \to 0}\frac{|\Delta x|}{\sqrt{2}|\Delta x|} = \frac{1}{\sqrt{2}} \neq 0.
    $$
    因此函数在原点处不可微。

    Q.E.D.
\end{pf}


17. 证明函数
$$
    f(x,y) =
    \begin{cases}
        (x^2 + y^2)\sin \dfrac{1}{\sqrt{x^2 + y^2}}, & x^2 + y^2 \neq 0, \\
        0,                                           & x^2 + y^2 = 0
    \end{cases}
$$
在点 $(0,0)$ 连续且偏导数存在，但偏导数在点 $(0,0)$ 不连续，而 $f$ 在原点 $(0,0)$ 可微。
\begin{pf}
    \begin{enumerate}
        \item 连续性：由于
              $$
                  \lim_{(x,y) \to (0,0)} f(x,y) = \lim_{\sqrt{x^2+y^2} \to 0} (x^2 + y^2)\sin\frac{1}{\sqrt{x^2 + y^2}} = f(0,0),
              $$
              因此函数在原点 $(0,0)$ 连续。

        \item 偏导数存在性：
              $$
                  f_x(0,0) = \lim_{\Delta x \to 0}\frac{f(\Delta x,0) - f(0,0)}{\Delta x} = \lim_{\Delta x \to 0}\frac{(\Delta x)^2\sin\frac{1}{|\Delta x|}}{\Delta x} = \lim_{\Delta x \to 0}\Delta x\sin\frac{1}{|\Delta x|} = 0.
              $$
              同理可得 $f_y(0,0) = 0$，因此偏导数在 $(0,0)$ 处存在。

        \item 偏导数不连续性：当 $\sqrt{x^2+y^2} \neq 0$ 时，
              $$
                  f_x(x,y) = 2x\sin\frac{1}{\sqrt{x^2 + y^2}} - \frac{x}{\sqrt{x^2+y^2}}\cos\frac{1}{\sqrt{x^2+y^2}}.
              $$
              令 $y = 0$，则 $\sqrt{x^2+y^2} = |x|$，于是
              $$
                  f_x(x,0) = 2x\sin\frac{1}{|x|} - \frac{x}{|x|}\cos\frac{1}{|x|}.
              $$
              由于 $\lim\limits_{x \to 0}\cos\frac{1}{|x|}$ 无限振荡，极限不存在，因此偏导数在点 $(0,0)$ 不连续。

              同理可证偏导数 $f_y(x,y)$ 在点 $(0,0)$ 不连续。

        \item 可微性：即证
              $$
                  \Delta z = f_x(0,0)\Delta x + f_y(0,0)\Delta y + o(h),
              $$
              其中 $h = \sqrt{(\Delta x)^2 + (\Delta y)^2}$。

              由已证 $f_x(0,0) = 0$, $f_y(0,0) = 0$，只需证明 $\Delta z = o(h)$。

              由于
              $$
                  \Delta z = f(\Delta x,\Delta y) - f(0,0) = h^2\sin\frac{1}{h},
              $$
              因此
              $$
                  \frac{\Delta z}{h} = h\sin\frac{1}{h} \to 0 \quad (h \to 0).
              $$
              即 $f$ 在原点 $(0,0)$ 可微。
    \end{enumerate}

    Q.E.D.
\end{pf}


18. 对下列函数，求出关于 $x$ 和 $y$ 的所有一阶和二阶（含混合）偏导数。

(1) $z = u\ln v$，其中 $u = x^2$，$v = \dfrac{1}{1+y}$;
\begin{sol}
    化简得
    $$
        z = x^2 \ln\frac{1}{1+y} = -x^2 \ln(1+y).
    $$
    因此一阶偏导数为
    $$
        \frac{\partial z}{\partial x} = -2x\ln(1+y), \quad \frac{\partial z}{\partial y} = -\frac{x^2}{1+y}.
    $$
    从而二阶偏导数为
    $$
        \frac{\partial^2 z}{\partial x^2} = \frac{\partial}{\partial x}\left(-2x\ln(1+y)\right) = -2\ln(1+y),
    $$
    $$
        \frac{\partial^2 z}{\partial x\partial y} = \frac{\partial}{\partial y}\left(-2x\ln(1+y)\right) = -\frac{2x}{1+y},
    $$
    $$
        \frac{\partial^2 z}{\partial y^2} = \frac{\partial}{\partial y}\left(-\frac{x^2}{1+y}\right) = \frac{x^2}{(1+y)^2}.
    $$
\end{sol}

(2) $z = u\arctan v$，其中 $u = \dfrac{xy}{x-y}$，$v = x^2y + y - x$.
\begin{sol}
    首先计算 $z$ 关于中间变量的偏导数：
    $$
        \frac{\partial z}{\partial u} = \arctan v, \quad \frac{\partial z}{\partial v} = \frac{u}{1+v^2}.
    $$
    其次计算 $u$, $v$ 关于 $x$, $y$ 的一阶偏导数：
    $$
        \frac{\partial u}{\partial x} = -\frac{y^2}{(x-y)^2}, \quad \frac{\partial u}{\partial y} = \frac{x^2}{(x-y)^2},
    $$
    $$
        \frac{\partial v}{\partial x} = 2xy - 1, \quad \frac{\partial v}{\partial y} = x^2 + 1.
    $$
    应用链式法则，得一阶偏导数：
    $$
        \frac{\partial z}{\partial x} = \frac{\partial z}{\partial u}\frac{\partial u}{\partial x} + \frac{\partial z}{\partial v}\frac{\partial v}{\partial x} = -\frac{y^2}{(x-y)^2}\arctan v + \frac{(2xy-1)u}{1+v^2},
    $$
    $$
        \frac{\partial z}{\partial y} = \frac{\partial z}{\partial u}\frac{\partial u}{\partial y} + \frac{\partial z}{\partial v}\frac{\partial v}{\partial y} = \frac{x^2}{(x-y)^2}\arctan v + \frac{(x^2+1)u}{1+v^2}.
    $$
    进而计算二阶偏导数：
    \begin{align*}
        \frac{\partial^2 z}{\partial x^2}
         & = \frac{\partial}{\partial x}\left(-\frac{y^2}{(x-y)^2}\arctan v\right)
        + \frac{\partial}{\partial x}\left(\frac{(2xy-1)u}{1+v^2}\right)                                                                                                                                       \\
         & =  \frac{2y^2}{(x-y)^3}\arctan v
        - \frac{y^2}{(x-y)^2}\cdot\frac{1}{1+v^2}(2xy-1)                                                                                                                                                       \\
         & \quad + \frac{(2y)(1+v^2)u + (2xy-1)(1+v^2)\frac{\partial u}{\partial x} - (2xy-1)u\cdot 2v\frac{\partial v}{\partial x}}{(1+v^2)^2}                                                                \\
         & = \frac{2y^2}{(x-y)^3}\arctan (x^2y + y - x)
        - \frac{y^2}{(x-y)^2}\cdot\frac{1}{1+(x^2y + y - x)^2}(2xy-1)                                                                                                                                          \\
         & \quad + \frac{\substack{(2y)(1+(x^2y + y - x)^2)(\frac{xy}{x-y})+ (2xy-1)(1+(x^2y + y - x)^2)\frac{-y^2}{(x-y)^2}- (2xy-1)(\frac{xy}{x-y})\cdot 2(x^2y + y - x)(2xy - 1)}}{(1+(x^2y + y - x)^2)^2}.
    \end{align*}
    \begin{align*}
        \frac{\partial^2 z}{\partial y^2}
         & = \frac{\partial}{\partial y}\left(\frac{x^2}{(x-y)^2}\arctan v\right)
        + \frac{\partial}{\partial y}\left(\frac{(x^2+1)u}{1+v^2}\right)                                                                                             \\
         & = -\frac{2x^2}{(x-y)^3}\arctan v
        + \frac{x^2}{(x-y)^2}\cdot\frac{1}{1+v^2}(x^2+1)                                                                                                             \\
         & \quad + \frac{(x^2+1)(1+v^2)\frac{\partial u}{\partial y} - (x^2+1)u\cdot 2v\frac{\partial v}{\partial y}}{(1+v^2)^2}                                     \\
         & = -\frac{2x^2}{(x-y)^3}\arctan (x^2y + y - x)
        + \frac{x^2}{(x-y)^2}\cdot\frac{1}{1+(x^2y + y - x)^2}(x^2+1)                                                                                                \\
         & \quad + \frac{\substack{(x^2+1)(1+(x^2y + y - x)^2)\frac{x^2}{(x-y)^2} - (x^2+1)(\frac{xy}{x-y})\cdot 2(x^2y + y - x)(x^2 + 1)}}{(1+(x^2y + y - x)^2)^2}.
    \end{align*}
    \begin{align*}
        \frac{\partial^2 z}{\partial x\partial y}
         & = \frac{\partial}{\partial y}\left(-\frac{y^2}{(x-y)^2}\arctan v\right)
        + \frac{\partial}{\partial y}\left(\frac{(2xy-1)u}{1+v^2}\right)                                                                                                                                        \\
         & = -\frac{2y(x-y)^2 + 2y^2(x-y)}{(x-y)^4}\arctan v
        - \frac{y^2}{(x-y)^2}\cdot\frac{1}{1+v^2}(x^2+1)                                                                                                                                                        \\
         & \quad + \frac{(2x)(1+v^2)u + (2xy-1)(1+v^2)\frac{\partial u}{\partial y} - (2xy-1)u\cdot 2v\frac{\partial v}{\partial y}}{(1+v^2)^2}                                                                 \\
         & = -\frac{2y(x-y)^2 + 2y^2(x-y)}{(x-y)^4}\arctan (x^2y + y - x)
        - \frac{y^2}{(x-y)^2}\cdot\frac{1}{1+(x^2y + y - x)^2}(x^2+1)                                                                                                                                           \\
         & \quad + \frac{\substack{(2x)(1+(x^2y + y - x)^2)(\frac{xy}{x-y}) + (2xy-1)(1+(x^2y + y - x)^2)\frac{x^2}{(x-y)^2} - (2xy-1)(\frac{xy}{x-y})\cdot 2(x^2y + y - x)(x^2 + 1)}}{(1+(x^2y + y - x)^2)^2}.
    \end{align*}
\end{sol}
\end{document}